% Raw Inscription Texts
\chapter*{Appendix B: Representative Inscriptions}
\addcontentsline{toc}{chapter}{Appendix B: Representative Inscriptions}

\section*{Sample Inscriptions with Operational Analysis}

This appendix provides representative inscriptions from the corpus with detailed operational breakdowns.

\subsection*{B.1: Simple Patterns}

\subsubsection*{Inscription H-001 (Harappa)}
\textbf{Object}: Square seal \\
\textbf{Sequence}: M361-M328-M411 \\
\textbf{Transliteration}: FISH-STROKE-TERMINAL \\
\textbf{Operational}: INST-COUNT(1)-CLOSE

\textbf{Analysis}:
\begin{itemize}
    \item M361 (Fish): Instantiation operator, introduces institutional context
    \item M328 (Stroke): Counting operator, cardinal value = 1
    \item M411 (Enclosure): Closure operator, validates transaction
\end{itemize}

\textbf{Interpretation}: Validated institutional record of single unit (commodity, transaction, or allocation).

\subsubsection*{Inscription M-045 (Mohenjo-daro)}
\textbf{Object}: Rectangular seal \\
\textbf{Sequence}: M346-M329-M329 \\
\textbf{Transliteration}: JAR-DOUBLE\_STROKE-DOUBLE\_STROKE \\
\textbf{Operational}: INST-COUNT(2)-COUNT(2)

\textbf{Analysis}:
\begin{itemize}
    \item M346 (Jar): Instantiation operator, container entity
    \item M329 (Double stroke): First count = 2
    \item M329 (Double stroke): Second count = 2 (total = 4?)
\end{itemize}

\textbf{Interpretation}: Container inventory with compound counting (possibly 2+2 or 2 of type A, 2 of type B).

\subsection*{B.2: Complex Patterns}

\subsubsection*{Inscription H-089 (Harappa)}
\textbf{Object}: Rectangular tablet \\
\textbf{Sequence}: M361-M346-M328-M328-M328-M411 \\
\textbf{Transliteration}: FISH-JAR-STROKE-STROKE-STROKE-TERMINAL \\
\textbf{Operational}: INST-QUAL-COUNT(1)-COUNT(1)-COUNT(1)-CLOSE

\textbf{Analysis}:
\begin{itemize}
    \item M361 (Fish): Primary instantiation (institutional marker)
    \item M346 (Jar): Qualification (specifies commodity type = jar/container)
    \item M328 $\times$ 3: Three counting operations = 3 units
    \item M411: Validation closure
\end{itemize}

\textbf{Interpretation}: Institutionally validated record of 3 jar-type commodities.

\subsection*{B.3: Maximum Length Examples}

\subsubsection*{Inscription D-023 (Dholavira)}
\textbf{Object}: Large signboard \\
\textbf{Length}: 10 signs \\
\textbf{Sequence}: M361-M346-M391-M328-M328-M347-M329-M373-M411-M412 \\
\textbf{Operational}: Complex multi-entity transaction with dual validation

\textbf{Analysis}: This represents one of the longest known inscriptions, showing:
\begin{itemize}
    \item Multiple entity instantiations
    \item Mixed counting operations
    \item Qualifications and relations
    \item Dual terminal markers (rare, possibly official/institutional signature)
\end{itemize}

\subsection*{B.4: Site-Specific Patterns}

\subsubsection*{Harappa Pattern A}
Characteristic pattern at Harappa: FISH-JAR-STROKES-TERMINAL \\
Frequency: 87 occurrences \\
Interpretation: Standard administrative receipt format

\subsubsection*{Mohenjo-daro Pattern B}
Characteristic pattern: JAR-STROKES (no terminal) \\
Frequency: 134 occurrences \\
Interpretation: Informal inventory marks, non-validated records

\subsubsection*{Dholavira Pattern C}
Characteristic pattern: FISH-ANIMAL-TERMINAL \\
Frequency: 23 occurrences \\
Interpretation: Institutional authorization, no quantity specification

\section*{Statistical Summary}

\begin{itemize}
    \item Total unique sequences: $\sim$2,800 (many inscriptions repeat)
    \item Most common length: 4--5 signs
    \item Maximum observed length: 17 signs (extremely rare)
    \item Median length: 4.2 signs
    \item Closure rate: $\sim$42\% of inscriptions have terminal markers
\end{itemize}
